
\begin{proposition}[Every class VII normalization component has Hopf minimal model]\label{cdp:st:C.2.UN.19}
For every proper holomorphic degeneration of smooth complex
two-tori with smooth three-dimensional total space, each
central component whose normalization resolution is of class
VII has a Hopf surface as its minimal model. This holds for
arbitrary multiplicities, arbitrary central singularities,
and arbitrary Euler number. In particular no Inoue surface
and no minimal class VII surface with positive second Betti
number is a normalization minimal model of such a component.

The comparison proving this assertion is the cyclic cover
obtained from $t=s^{m_i}$ for that component. Its stabilizer
acts on the minimal model of the covering component, and
the quotient is bimeromorphic to the original normalization.
\end{proposition}

\begin{proof}
We first identify the quotient in the actual base change.
Normalize $U\times_{\Delta_t}\Delta_s$, where
$t=s^{m_i}$. Its cyclic group $\mu_{m_i}$ acts over $U$.
Choose a resolution of the total space and central divisor
equivariant for this finite group. Canonical analytic
resolution has this property: the invariants defining its
centres are preserved by every local isomorphism, so the
action lifts through every blowup
\cite[\S1, universal and canonical desingularization]{BierstoneMilman1997}.
The algorithm is an isomorphism over the punctured tube
and at the generic smooth points used in (C.SC.2).

Let $E$ be a chosen smooth strict component above $S_i$,
and let $G\subset\mu_{m_i}$ be its stabilizer. At a
generic point of $S_i$ the equations are exactly
\[
                 w^{m_i}=u,\qquad z=s/w.
 \tag{C.HQ.1}\label{cdp:eq:C.HQ.1}
\]
Consequently the reduced inverse image over a dense
smooth open subset of $S_i$ is a $\mu_{m_i}$-torsor.
The cyclic group permutes its connected components
transitively. The stabilizer of a selected component
acts freely and transitively on each fibre of that
component. Equivalently, the degree of $E\to S_i$
at the generic point is $|G|$, and quotienting by
$G$ gives degree one. The proper map to normalization
therefore descends to a proper bimeromorphic map
\[
                         E/G\longrightarrow S_i^\nu.
 \tag{C.HQ.2}\label{cdp:eq:C.HQ.2}
\]
This argument remains valid when the root cover is
disconnected over $S_i$; the stabilizer, rather than
the whole base-change group, is the group in \eqref{cdp:eq:C.HQ.2}.

Write $Z_i$ for a smooth normalization resolution and
$M_i$ for its minimal model. The case $b_2(M_i)=0$
already has Hopf minimal model by \ref{cdp:st:C.2.UN.17}. Every
minimal class VII surface with $b_2>0$ has algebraic
dimension zero \cite[p.~256]{DOT2000}. It therefore
suffices to treat $a(M_i)=0$, by an argument which
also covers the zero-Betti case. The field comparison
in (C.CF.9) gives
$a(E)=0$, and \ref{cdp:st:C.2.UN.16} says that the unique smooth
minimal model $E_0$ is a torus or a Hopf surface.

Every element of $G$ induces a birational automorphism
of $E_0$. On a nonruled minimal compact surface this
is holomorphic, by uniqueness of the minimal model.
Thus the contraction $\rho:E\to E_0$ is equivariant:
the identity $\rho g=g_0\rho$ holds on the common
open isomorphism and hence everywhere. It descends
to a proper bimeromorphic map
\[
                    E/G\longrightarrow E_0/G.
 \tag{C.HQ.3}\label{cdp:eq:C.HQ.3}
\]
The actions on $E_0$ are faithful after removing
their kernel, which does not change the quotient.
Resolving the graph of \eqref{cdp:eq:C.HQ.2}--\eqref{cdp:eq:C.HQ.3} gives
actual proper bimeromorphic comparisons between
$Z_i$ and a resolution of $E_0/G$.

We exclude the torus alternative by its first
cohomology. For a finite holomorphic quotient
$q:A\to A/G$, proper base change gives
$R^jq_*\mathbb Q=0$ for $j>0$, and averaging gives
$(q_*\mathbb Q)^G=\mathbb Q$. Taking invariants is
exact over $\mathbb Q$, so the actual pullback map
identifies
\[
                     H^1(A/G,\mathbb Q)
                          \simeq H^1(A,\mathbb Q)^G.
 \tag{C.HQ.4}\label{cdp:eq:C.HQ.4}
\]
For a complex torus the holomorphic action preserves
$H^{1,0}\oplus H^{0,1}$ and conjugation interchanges
the two invariant summands. Hence the dimension in
\eqref{cdp:eq:C.HQ.4} is even. Each singularity of a cyclic surface
quotient has a resolution with a connected chain
of rational curves: after linearization, the
quotient cone is the first quadrant in
$\mathbb Z^2+\frac1n(a,b)\mathbb Z$, and a regular
fan subdivision gives that resolution. Removing
pseudoreflections first only replaces the local
coordinates by their invariant powers. Thus every
exceptional fibre has $H^1=0$ and is connected.
Proper base change and the degree-one part of
Leray show that resolving $A/G$ leaves $H^1$
unchanged. Birational maps between smooth compact
surfaces preserve $H^1$, as follows from their
factorization into point blowups. The comparisons
\eqref{cdp:eq:C.HQ.2}--\eqref{cdp:eq:C.HQ.3} would therefore make $b_1(Z_i)$
even, contradicting $b_1(Z_i)=1$.

It remains to prove that the cyclic quotient of
the Hopf surface $H=E_0$ is again Hopf. We give
the fixed-point and quotient argument explicitly.
Let
\[
            H=(\mathbb C^2\setminus\{0\})/\Gamma.
 \tag{C.HQ.5}\label{cdp:eq:C.HQ.5}
\]
The deck group contains an infinite cyclic
subgroup of finite index generated by a
holomorphic contraction. Equivalently $H$ has
a finite unramified primary Hopf cover; the
normal forms and the finite-cover description
are recalled in \cite[\S1.1]{Nakamura2008VII}
and \cite[Theorem~2.1]{Khaled2024Hopf}.
Every deck transformation extends across the
origin by Hartogs, as does its inverse. These
extensions fix the origin and are biholomorphic.
The same is true for any lift $\widetilde g$
of a holomorphic automorphism $g$ of $H$.

Define the real character
\[
 \delta:\Gamma\longrightarrow\mathbb R,
 \qquad \delta(\gamma)=\log|\det(d\gamma)_0|.
 \tag{C.HQ.6}\label{cdp:eq:C.HQ.6}
\]
It is nonzero, since its value on a contraction
is negative. For any automorphism lift,
conjugacy of derivatives gives
\[
 \delta(\widetilde g\gamma\widetilde g^{-1})
                         =\delta(\gamma).
 \tag{C.HQ.7}\label{cdp:eq:C.HQ.7}
\]
A primary Hopf cover has rational Betti numbers
$(1,1,0,1,1)$: this follows directly from the
covering spectral sequence for
$\mathbb C^2\setminus\{0\}\simeq S^3$ and the
infinite cyclic deck group, whose action on
$H^3(S^3,\mathbb Q)$ preserves orientation.
Transfer from the finite primary cover gives
$b_1(H)\le1$ and $b_2(H)=0$, while \eqref{cdp:eq:C.HQ.6}
gives $b_1(H)\ge1$. Hence $b_1(H)=1$.
Equation \eqref{cdp:eq:C.HQ.7} says that every holomorphic
automorphism acts identically on $H^1(H,\mathbb R)$.
Poincar\'e duality and its complex orientation
then give the identity action on $H^3$ as well.
Its Lefschetz number is therefore
\[
                 L(g)=1-1+0-1+1=0.
 \tag{C.HQ.8}\label{cdp:eq:C.HQ.8}
\]

For a nonidentity finite-order automorphism,
the fixed locus is a disjoint union of smooth
compact curves and isolated points, by finite
holomorphic linearization. Each fixed curve $C$
is smooth. Since $H^2(H,\mathbb Q)=0$, one has
$C^2=K_H\cdot C=0$. The actual adjunction
isomorphism $K_C=(K_H\otimes\mathcal O_H(C))|_C$
therefore gives $2g(C)-2=0$. Thus all the
curve components of the fixed locus have Euler
number zero. For a finite-order
map the topological Lefschetz number equals
the Euler number of its fixed locus. Indeed,
use an equivariant triangulation and barycentric
subdivision; each invariant simplex is then
fixed pointwise, so the alternating trace
on simplicial cochains counts exactly the
fixed simplices. Equation \eqref{cdp:eq:C.HQ.8} consequently
shows that there are no isolated fixed points.

At any point of $H$, the stabilizer in the
effective cyclic group $G$ is cyclic. Its
generator has a fixed curve through that
point and, in linearizing coordinates, acts
by
\[
                (x,y)\longmapsto(x,\zeta y),
 \qquad
 \mathbb C\{x,y\}^{\langle\zeta\rangle}
                    =\mathbb C\{x,y^{|G_x|}\}.
 \tag{C.HQ.9}\label{cdp:eq:C.HQ.9}
\]
Here $G_x$ denotes the stabilizer, and
$\zeta$ has its full order; a finite-order
germ with identity derivative is the identity.
This proves that the quotient $Q=H/G$ is
smooth, with the displayed local quotient
map even at its branch curves.

The invariant-cohomology comparison in
\eqref{cdp:eq:C.HQ.4}, now in every degree and combined
with \eqref{cdp:eq:C.HQ.7}, gives
\[
                         b_1(Q)=1,
                         \qquad b_2(Q)=0.
 \tag{C.HQ.10}\label{cdp:eq:C.HQ.10}
\]
The finite map $H\to Q$ injects every space
of pluricanonical sections by the actual
Jacobian pullback. Since a Hopf surface
has $\kappa=-\infty$, this proves
$\kappa(Q)=-\infty$. Also $Q$ is minimal:
an exceptional curve of self-intersection
$-1$ would have a nonzero real $H^2$ class,
contrary to \eqref{cdp:eq:C.HQ.10}.

There is a curve on $Q$. Explicitly, the
primary Hopf normal form
$(z_1,z_2)\mapsto(\alpha z_1+a z_2^r,\beta z_2)$
contains the elliptic curve given by $z_2=0$;
its finite image in $H$, and then in $Q$,
is a nonempty curve. A positive multiple
of that curve is homologically trivial by
\eqref{cdp:eq:C.HQ.10}. The divisor classification
\cite[Remark~2.7]{DlousskyTeleman2012} applied
to the minimal class VII surface $Q$ now
gives a Hopf surface: its Enoki alternative
has positive second Betti number and is
excluded by \eqref{cdp:eq:C.HQ.10}. We have proved the
cyclic-quotient assertion by actual quotient
maps, including their branch points.

Finally \eqref{cdp:eq:C.HQ.2}--\eqref{cdp:eq:C.HQ.3} make $Z_i$
bimeromorphic to this smooth Hopf surface
$Q$. Uniqueness of the class VII minimal
model proves that $M_i\simeq Q$. This
proves the assertion for every component
in the original, unaltered degeneration.
\end{proof}


